\documentclass[11pt]{article}
\usepackage[margin=1in]{geometry}
\usepackage{enumitem}
\usepackage{nopageno}

\begin{document}

\noindent \textbf{Facilities, Equipment, and Other Resources}

\begin{itemize}

\item \textbf{PSU College of Engineering Institutional Strategic Alignment}
Penn State University is a land-grant institution deeply committed to providing access to life-changing educational opportunities, particularly for underrepresented populations. The College of Engineering is fully committed to the success and long-term impact of the Phase I CAMEL Collaborative Network (CAMEL-CN). This project aligns with strategic institutional goals of advancing pedagogical innovation and building a nationally competitive STEM workforce.

\item \textbf{PSU College of Education Institutional Strategic Alignment}
The Penn State College of Education is nationally identified as having major research and training programs as well as outreach initiatives focused on the enhancement of STEM teaching and learning. The College of Education’s Grants and Contracts Office provides post-award administration in the form of budgetary oversight, implementation of research methodologies, and mechanisms to ensure on-time reporting to funding agencies. The use of technological tools to enhance teaching, learning, and research is a major goal of Penn State University, making available knowledgeable partners for multimedia and other technologically enhanced projects.

\item \textbf{Office for Digital Learning (ODL)}
Reporting to the Associate Dean for Education, the ODL serves as a core institutional resource for promoting innovative engineering education both on campus and online. The CAMEL-CN network will leverage ODL’s specialized infrastructure to ensure the pedagogical integrity and technical accessibility of the project's digital outputs.
    \begin{itemize}[noitemsep]
        \item \textbf{Instructional Design and Technology:} ODL maintains a curated collection of teaching strategies and educational technology tools specifically designed for engineering contexts. These resources will support the research team in developing effective digital interfaces for the Student Math Modeling Learning Dataset and ensuring that process-level telemetry is captured seamlessly within the learning environment.
        \item \textbf{Pedagogical Support Artifacts:} The project will utilize ODL’s extensive library of tutorials and templates to standardize the design of the mathematical modeling modules. This institutional resource ensures that the instructional materials produced by the network are grounded in evidence-based teaching practices and are accessible to high school educators without requiring advanced programming skills.
    \end{itemize}

\item \textbf{Center for Science and the Schools (CSATS)}
CSATS serves as the project’s primary operational bridge between STEM research and the K-12 community. The Center provides a robust administrative and pedagogical infrastructure to ensure the effective co-design and implementation of the proposed activities.

    \begin{itemize}[noitemsep]
        \item \textbf{Program Development and Activity Design:} CSATS faculty, including \textbf{Amber Cesare} (Instructional Design Specialist), provide expertise in education research-based practices. This operational unit works closely with the science and engineering faculty to develop activities that maximize student learning by building knowledge of both industrial content and research practices.
        
        \item \textbf{Recruitment and Marketing Infrastructure:} CSATS maintains an extensive communication network with Pennsylvania school districts, including dedicated listservs and strong connections to the rural education community. These resources facilitate the recruitment of the seven committed rural and Title I districts and the planned expansion to 4,500 students.
        
        \item \textbf{Participant Registration and Management:} CSATS provides a proprietary website system to manage the registration of all program participants. Operationally, the Center handles all necessary preparations and arrangements to provide a positive participant experience, including the coordination of travel, subsistence, and the disbursement of project-funded stipends and supplies under the direction of \textbf{Kathleen Hill} (Director).
        
        \item \textbf{Program Evaluation Infrastructure:} CSATS provides the specialized resources required to assess the success of the project’s broader impacts. This infrastructure supports the identification of measurable outcomes, the development of specialized evaluation instruments, and the rigorous collection and analysis of data required for reporting findings to the NSF.
        
        \item \textbf{AI/Computational Visualization Support:} The Center will provide operational space and oversight for a \textbf{Computer Science/AI Education Specialist}. This specialist will utilize CSATS’s technical resources to develop the computational tools and visualization interfaces that allow K-12 teachers to navigate authentic industrial datasets without requiring advanced programming skills.

        \item \textbf{Science Education Wing and Resource Center: } Dr. Hill and other CSATS employees and graduate students have access to newly-renovated science education classrooms with teaching and laboratory space. The Science Ed Resource Center houses science teaching materials including a wide variety of technologies, kits, and books.

        \item \textbf{Krause Innovation Studio and Pedagogy Innovation Lab:} Dr. Hill and other CSATS employees and graduate students have access to the Krause Innovation Studio which is at the forefront of supporting research in teaching and learning with emerging and innovative digital tools, including data visualization. All video recording will be conducted using the IRIS Connect system. IRIS Connect provides secure video recording through an integrated iOS app and web-based platform, which enables ongoing feedback and collaboration between teacher and researcher. With the assistance of the Krause Innovation Studio, the IRIS Connect systems will be used to collect of observational data from classrooms. The video recordings will also be used to facilitate professional development. The Krauss Pedagogy Innovation Lab is a 1,400-square-foot space adjacent to the existing Krause Studios, designed to introduce our teachers to new and emerging technologies and approaches to teaching and engaging their students. It is a vibrant and versatile classroom that supports the discovery of new pedagogies through making and other creative activities. The space includes the kinds of equipment common to makerspaces, such as 3D printers, a laser cutter and woodworking equipment that allow inservice and preservice teachers to design and build real objects, taking an idea from concept to physical reality all in one space. A one-button podcast studio provides opportunities to create within the audio and video realm. The PIL also includes several elements to facilitate classroom teaching, presentations and collaboration, such as a whiteboard/glassboard, wall-mounted touch-screen monitors, and a retractable stage for spoken-word and other performance-based presentations.
 
    \end{itemize}

\item \textbf{Data Science Education Infrastructure (NSF HDR DSC)}
The project leverages the substantial intellectual and methodological resources established through the ongoing Harnessing the Data Revolution (NSF HDR DSC) initiative at Penn State. 
    \begin{itemize}[noitemsep]
        \item \textbf{Curricular Framework:} The HDR project has developed and validated an evidence-based blueprint for integrating data science with engineering curricula. This framework provides the CAMEL-CN project with systematic protocols for data translation.
        \item \textbf{Module Repository:} The network has access to an existing online repository of six data science courses tailored for engineering contexts. These materials serve as a structural precedent for the design of high school Algebra and Statistics modules.
        \item \textbf{Stakeholder Network:} The HDR project has established a community of stakeholders and industry partners who provide ongoing feedback on the industrial relevance of data modeling tasks.
    \end{itemize}
    
\item \textbf{National User Facility: 2D Crystal Consortium (2DCC-MIP)}
Co-PI Reinhart will leverage the resources of the 2DCC-MIP, an NSF-funded National Materials Innovation Platform.
\begin{itemize}[noitemsep]
    \item \textbf{Industrial Data Repositories:} The project will have full access to the 2DCC-MIP Lifetime Sample Tracking (LiST) database, which contains over 20,000 specialized material samples and associated characterization data. 
    \item \textbf{Composable Courseware Infrastructure:} The project will utilize the existing "Composable Courseware" (CC) architecture and authorship layer, which provides the software foundation for the modular construction of plug-and-play educational modules.
\end{itemize}

\item \textbf{Building Energy and AI Infrastructure Systems}
Co-PI Zuo provides access to high-dimensional time-series datasets from several NSF-funded cyber-physical research initiatives. These resources provide students with authentic engineering challenges involving sensor noise, high dimensionality, and provenance artifacts.
    \begin{itemize}[noitemsep]
        \item \textbf{District Cooling and Heating Networks:} Access to operational data from campus-scale systems (e.g., CU Boulder), including sensor records of temperatures, flow rates, and electricity consumption. These data include authentic ``messiness'' such as physically implausible readings under low-flow conditions and sparse metering in aging infrastructure.
        \item \textbf{Data Center Cooling System:} High-frequency time-series data from actual university data centers (e.g., UMass), featuring abrupt control-driven mode changes and strong coupling between weather, IT load, and cooling response.
        \item \textbf{Virtual Testbed for Net-Zero Energy Communities (NZEC):} Access to an open-source virtual testbed based on real-world communities, consisting of residential, office, and retail buildings with integrated solar PV and HVAC grid interactions.
    \end{itemize}

\item \textbf{Materials Informatics Training Infrastructure (NSF NRT)}
The network will leverage the pedagogical and computational infrastructure of the ongoing NSF National Research Traineeship (NRT) at Penn State. 
    \begin{itemize}[noitemsep]
        \item \textbf{Proven Adaptation Protocols:} Led by Co-PI Reinhart, the NRT has developed systematic protocols for adapting materials science research datasets for data science training. This existing framework will serve as a foundational resource for the CAMEL-CN Data Translation Framework (Task 3).
        \item \textbf{Interdisciplinary Training Environment:} The NRT provides a model for cross-disciplinary communication between data originators and learners. This established environment ensures that the network has access to refined strategies for presenting ``messy'' materials informatics data to non-expert audiences.
    \end{itemize}

\item \textbf{Institute for Computational and Data Sciences (ICDS-ACI)}
The project will generate tens of gigabytes of raw process-level telemetry and tens of millions of annotated events. This data will be hosted through the ICDS-ACI systems at Penn State.
\begin{itemize}[noitemsep]
    \item \textbf{Active Storage:} High-resolution telemetry and process data will be managed via DDN 12KX40 and GS7K flash storage array systems, ensuring the low-latency access required for real-time teacher dashboards and the FastAPI backend.
    \item \textbf{Archival and Preservation:} Long-term archival services utilize an Oracle SL8500 Tape Library and FS1 flash storage appliance. Multiple copies of archived data are managed through Oracle Hierarchical Storage Manager (HSM) to safeguard against corruption.
    \item \textbf{Security Infrastructure:} ICDS enforces "Least Privilege" access, strong encryption in-flight and at rest, and multi-factor authentication, ensuring that Restricted Tier demographic data remains protected and compliant with IRB requirements.
\end{itemize}

\item \textbf{Center for Socially Responsible Artificial Intelligence (CSRAI)}
PI Napolitano is an affiliate member of CSRAI. The network will leverage the Center’s expertise in the ethical implications of AI deployment. CSRAI resources will support the project in developing frameworks for non-evaluative student monitoring and responsible real-time telemetry intervention, ensuring student agency and privacy are preserved.

\item \textbf{Research Data Stewardship and Curation}
The project will benefit from the Penn State Research Data Stewardship Program, a joint initiative between the University Libraries and the Office of the Senior Vice President for Research. This program provides institutional infrastructure to ensure compliance with University Policy RP15 (Research Data Management) and federal mandates for data accessibility and reproducibility.
    \begin{itemize}[noitemsep]
        \item \textbf{Data Management and Consulting:} Led by the Director of Research Data Stewardship, this unit provides specialized training and consulting for data management planning, statistical analysis, and implementation. This ensures that the project’s high-dimensional datasets are curated for long-term usability.
        \item \textbf{Institutional Repository (ScholarSphere):} The project has access to ScholarSphere, a secure repository maintained by the University Libraries Strategic Technologies Department. ScholarSphere provides a platform for the open sharing of research articles, industrial datasets, and creative works produced by the CAMEL-CN network.
        \item \textbf{Data Curation Network (DCN):} Through Penn State’s membership in the DCN, the project leverages a national network of 23 institutional repositories. This provides the research team with access to expert cross-institutional curation, ensuring that data is findable, accessible, interoperable, and reusable (FAIR).
        \item \textbf{Quality Assurance and Compliance:} The Office for Research Protections provides dedicated Quality Assurance teams to conduct data management reviews. This infrastructure supports the project in maintaining secure storage on University-approved devices and managing data retention for the required minimum of seven years post-publication.
    \end{itemize}

\item \textbf{Laboratory and Office Space}
The network utilizes dedicated computational and research space within the colleges of Education, Engineering and Earth and Mineral Sciences at Penn State.
\begin{itemize}[noitemsep]
    \item \textbf{Computational Facilities:} The University Park campus provides high-speed networking and secure server infrastructure capable of hosting the FastAPI backend and the Composable Courseware authorship tools.
    \item \textbf{Collaborative Workspace:} PI's / CoPI's lab provides standard office equipment and collaborative software tools for the co-production of expert-coded annotations by the learning science and data science teams.
\end{itemize}

\item \textbf{Graduate and Professional Support}
Graduate students supported by this project have access to the Penn State Graduate Career Services and the Graduate Writing Center. These entities provide career counseling, training in proposal preparation, and guidance on professional practices, which will be integrated into the project's Mentoring Plan.

\end{itemize}

\end{document}